\section{NestSection: Algorithm Recap}
\label{sec:nestsection}

We briefly recap the NestSection algorithm from companion paper B.13
\citep{dellaLibera2026nestSection}. Readers familiar with B.13 may proceed
directly to Section~\ref{sec:analysis}.

\subsection{The Spine Construction}
\label{sec:nestsection:spine}

\begin{definition}[Spine]
  Given house count $H$ and senate count $S$ with $g = \gcd(H, S) > 1$,
  the \emph{spine} is a $g$-way partition of the state into $g$ geographic
  super-regions, each containing exactly $H/g$ house districts and $S/g$
  senate districts.
\end{definition}

The spine is computed by METIS $g$-way partitioning of the state's
census-unit graph with equal population targets (each super-region receives
$T/(g)$ residents). Within each spine super-region:
\begin{itemize}
  \item The house sub-map is the $(H/g)$-way partition of the super-region.
  \item The senate sub-map is the $(S/g)$-way partition of the same super-region.
  \item Nesting is guaranteed: each senate district within the super-region
    is the union of $(H/g) / (S/g) = H/S$ house districts.
\end{itemize}

This construction guarantees nesting by construction: house districts are drawn
as leaves of the spine tree, and senate districts are drawn by grouping
consecutive house districts within each spine super-region.

\subsection{The Ratio Condition}
\label{sec:nestsection:ratio}

NestSection requires $g = \gcd(H, S) > 1$. When $g = 1$, $H$ and $S$ are
coprime: there is no non-trivial common divisor, and the spine degenerates to
a single region containing all districts. This does not provide any of the
compactness benefits of spine construction, and nesting cannot be guaranteed
by the algorithm.

\begin{proposition}[Spine triviality when $g=1$]
  If $\gcd(H, S) = 1$, the NestSection spine has $g = 1$ super-region
  (the whole state). The house map is the $H$-way partition of the whole state
  and the senate map is the $S$-way partition of the whole state; no nesting
  constraint can be derived from the spine.
\end{proposition}

For $g = 1$ states, the constitutional nesting requirement (where it exists)
cannot be satisfied by NestSection. These states require either an ad-hoc
nesting procedure or must use a different chamber ratio.

\subsection{NestSection CLI Usage}
\label{sec:nestsection:cli}

\begin{verbatim}
# Washington: H=98, S=49, gcd=49 (ratio 2:1)
redist state --state WA --chamber nest \
             --house-districts 98 --senate-districts 49 \
             --resolution block_group --year 2020

# Minnesota: H=134, S=67, gcd=67 (ratio 2:1)
redist state --state MN --chamber nest \
             --house-districts 134 --senate-districts 67 \
             --resolution tract --year 2020
\end{verbatim}

The \texttt{--chamber nest} flag triggers NestSection mode. The spine count
$g = \gcd(H, S)$ is computed automatically; the spine super-regions are
partitioned first, then each super-region is independently redistricted for
house and senate.
